% Transcription — fonds Alexandre Grothendieck, shelfmark 120, batch 1 (pages 1-20).
% Established from the facsimile archives/batches/120/batch-01.pdf (a copy of
% 120.pdf fetched through the project's relay,
% https://grothendieck-relay.onrender.com/source/120.pdf; PDF page k is the
% archivists' page k, no cover sheet), per the skill transcribe-grothendieck.
% The folder is twenty-four pages; this is the first of two batches (pages
% 21-24 are batch 2, transcribed separately; its notation — R, R_0,
% \mathcal{M}, the labels (M1)-(M11), \smallsetminus, \pi_0 — is followed here).
%
% Pass: Opus 5.5 (claude-opus-5-5), 2026-09-26 — first pass, unchecked against the pages by a human.
%
% Pages skipped: none as unrelated. Page 1 is a cover (pink sheet) inscribed
% in ink « Topologie modérée » (the first word written over a struck word,
% perhaps « Structures », the final -es of « modérées » corrected), with a
% very faint pencil « 1978 ... (Autour de l'Éloge) » below; it takes no
% \page and its words become the section heading, with a \note.
% Pages summarised: none; but three long slanting additions squeezed between
% the lines (pp. 5, 6, 7) are only partly legible and are described in a
% \note with the fragments read, rather than set in full.
% Pages transcribed: 2-20, all his hand, one ink, fast drafting.
%
% His pagination (top right): 1 (p. 2), 2 (p. 4), [3, faded] (p. 6),
% 4 (p. 8; p. 7 carries none legible), 5 (p. 10), 6 (p. 12), 7 (p. 14), 8 (p. 16), 9 (p. 18), 10 (p. 20):
% he numbers roughly every second leaf. Batch 2's « 11 » on p. 22 continues
% this run exactly (p. 20 = 10, p. 22 = 11).
%
% What the batch is about. An axiomatic start on tame topology (« topologie
% modérée », his term, as in Esquisse d'un programme). (I) Modèles
% R-modérés: for a set R, classes M_n of subsets of R^n, axioms (M1)-(M4)
% (stability under the maps induced by [1,n] -> [1,m], inverse images by
% surjective ones, intersections with inverse images, products), their
% reduction to permutations, face and degeneracy maps; trivial cases
% excluded ((*), (**), (***)); R_0 = points x with {x} in M_1; examples.
% (II) Espaces modérés: tame maps (graph in M_{I+J}), the category of models
% and the equivalent category Esp_{M_*} with a faithful conservative
% « transportable » forgetful functor; finite products, final object and
% finite limits under (M5) R_0 != Ø; (M0) R in M_1; R as an object of
% Sigma = Esp_{M_*}, every object a subobject of some R^n, and a
% characterisation of M_* by (Sigma, R, X |-> |X|). (III) Finite colimits:
% the discrete tame structure on a finite set; a set-theoretic colimit
% carrying a tame structure making the g_i tame is the colimit in Sigma;
% a second (M5) (stability of M_n under finite unions), unions of tame
% parts, the amalgamated-sum question X_1 +_Z X_2 and a criterion by a
% finite family of pairs of functions to R separating points; quotients
% X/R_f by the equivalence relation of a tame map; (M6) (a tame Z in X is a
% fibre f^{-1}(e) of finitely many tame functions, OPS separating points of
% X \ Z, i.e. X/Z exists strictly) and (M7) (extension of tame functions);
% under M5-M7 amalgamated sums exist and commute with the forgetful functor;
% the embedding diagram (*) proving it. (IV) Ind-structures modérées
% strictes (ind-tame spaces), (M8) R the union of its tame parts, finite
% limits and sums, locally tame parts, quotients X/Y and amalgamated sums of
% ind-tame spaces. (V) Relation with topologies: (M9) R Hausdorff and tame
% parts compact, so Esp_{M_*} -> compact spaces; ind-tame spaces with the
% colimit topology; (M10) every point of R has a fundamental system of
% neighbourhoods in M_1; tameness is local, the sheaf of tame maps;
% « espaces loc. modérés ». (VI) Relation with a field structure on R:
% (M11) a) 0 != 1 in R_0, b) x-y and xy tame on A in M_2, c) inverses of
% A in R^* in M_1 lie in some B in M_1 (OPS B = A^{-1}); tame functions form
% a subring, R_0 a subring of R. Batch 2 picks up with ind-tame spaces and
% (M11) c) (the inverse map).
%
% Notation fixed once for the batch, following batch 2: plain R, R_0;
% \mathcal{M}_n, \mathcal{M}_I, \mathcal{M}_* for his script M;
% \mathfrak{P} for the power set, \mathfrak{S}_n for the symmetric group (and
% as he writes it, p. 2, where \mathcal{M}_n is meant: noted);
% \mathrm{Esp}_{\mathcal{M}_*} for his « Esp_{M*} », \Sigma for the same;
% \mathfrak{X}, \mathfrak{Y}, \mathfrak{Z} for his gothic letters on
% pp. 15-16 (ind-tame spaces); \mathcal{R}_f for the equivalence relation of
% f; the axioms as (M1) ... (M11), parentheses where he has them; his
% « M1 : M5 » (p. 14) kept with the colon, as batch 2 keeps « M1 : M7 ».
% Two labels are used twice by him: (M5) is R_0 != Ø on p. 6 and stability
% under finite unions on p. 9 (and « impliqué par M5 » p. 14 refers to
% the latter); noted, not regularised.
%
% For batch 2: the prefix « ind- » is settled by p. 14, where the section
% title and its interlinear variant write « Ind-structures », « Ind-espaces
% modérés » with an unmistakable capital I; the lower-case « ind-modéré »
% of pp. 14-19 is the same stroke as batch 2's. The outlined \mathbb{D}
% opening p. 21 is not explained by these pages: p. 20 ends mid-sentence
% (« En particulier R_0 est un sous-anneau de R, et les ... »).
%
% Readings to check first: the faint pencil on the cover; the struck words
% of the marginal title on p. 2; the slanting additions of pp. 5, 6, 7;
% the whole top of p. 8 (four struck lines); p. 11 first lines; the
% direction of the hooks in diagram (*) p. 13; the marginal notes of pp. 14,
% 19, 20; « dénombrable » p. 17.
% Apparatus: 130 \ill{}, 74 \uncertain{}.

\documentclass[11pt,a4paper]{article}
% Le préambule commun à toutes les transcriptions du fonds Grothendieck.
%
% Il tient en une idée : **l'incertitude se compose, elle ne se résout pas.**
% Une transcription qui lisse un mot douteux a détruit la seule chose pour
% laquelle on la faisait — un lecteur doit pouvoir savoir, à chaque endroit,
% ce qui était sur le feuillet et ce qui a été ajouté. Les six macros qui
% suivent sont donc l'appareil critique, et non de la décoration.
%
% Les mêmes commandes sont reconnues par `scripts/render.mjs`, qui produit la
% vue de lecture du site. Une macro ajoutée ici doit l'être aussi là-bas,
% faute de quoi le rendu échouera bruyamment — ce qui est voulu : un rendu
% silencieusement faux est pire qu'un rendu absent.

\ProvidesPackage{grothendieck}[2026/08/08 Transcriptions du fonds Grothendieck]

\usepackage{amsmath,amssymb,amsthm}
\usepackage{mathrsfs}
\usepackage{stmaryrd}
% Les diagrammes commutatifs sont partout dans ce fonds : c'est la langue
% même de la géométrie algébrique de Grothendieck, et une transcription qui
% les rendrait en prose ne transcrirait rien.
\usepackage{tikz-cd}
% Français par défaut — c'est la langue des feuillets — mais l'anglais est
% chargé aussi : la traduction et le résumé sont des documents anglais, et la
% typographie française (espaces avant les deux-points) y serait une faute.
% Ces documents basculent par \selectlanguage{english} en tête de corps.
\usepackage[english,french]{babel}
\usepackage{fontspec}
\usepackage{geometry}
\geometry{a4paper,margin=25mm,marginparwidth=28mm}
\usepackage{marginnote}
\usepackage{xcolor}
% ulem, not soul. `soul`'s \st and \ul are built on a character-by-character
% scan that XeLaTeX's font machinery breaks: the document fails with an
% undefined control sequence rather than degrading. ulem does the same two jobs
% and is Unicode-engine safe. `normalem` stops it hijacking \emph, which the
% transcriptions use for Grothendieck's own emphasis.
\usepackage[normalem]{ulem}
\usepackage{hyperref}
\hypersetup{colorlinks=true,linkcolor=black,urlcolor=black,pdfborder={0 0 0}}

% --- Métadonnées du lot -----------------------------------------------------
% Reprises telles quelles par le script de rendu, qui les affiche en tête de
% la vue de lecture. Elles fixent ce que le fichier prétend couvrir : sans
% elles, une transcription flotte sans point d'attache dans un fonds de seize
% mille feuillets.

\newcommand{\folder}[1]{\gdef\grothFolder{#1}}
\newcommand{\batch}[1]{\gdef\grothBatch{#1}}
\newcommand{\leaves}[2]{\gdef\grothFirst{#1}\gdef\grothLast{#2}}
\newcommand{\foldertitle}[1]{\gdef\grothFoldertitle{#1}}
\gdef\grothFolder{?}\gdef\grothBatch{?}\gdef\grothFirst{?}\gdef\grothLast{?}\gdef\grothFoldertitle{}

% --- Le résumé --------------------------------------------------------------
%
% Il ouvre la lecture modernisée, sous le titre « Résumé », et il est composé
% en petit. Ce n'est pas un ornement : c'est le seul passage du document qui
% ne soit pas des mathématiques, et un lecteur doit pouvoir le reconnaître
% avant de l'avoir lu — puis le sauter s'il connaît déjà le sujet, ou s'y
% arrêter s'il ne le connaît pas. Le filet qui le suit marque la reprise du
% corps.

\newenvironment{resume}
  {\small\setlength{\parskip}{0.45em}}
  {\par\medskip\hrule\medskip}

% --- Mots-clés ---------------------------------------------------------------
%
% En anglais, en clôture du résumé de la lecture modernisée : le vocabulaire
% moderne sous lequel chercher ce que les feuillets construisent. Le manifeste
% du site (`npm run manifest`) les extrait pour étiqueter la cote entière —
% cette ligne est la source unique des tags, il n'y a pas de fichier à côté.
\newcommand{\keywords}[1]{\par\smallskip\noindent
  {\footnotesize\textsc{Keywords} --- \textit{#1}}\par}

% --- L'appareil critique ----------------------------------------------------

% Le numéro de feuillet, en marge. C'est l'ancre qui permet de régler un
% désaccord sur une formule en regardant *un* feuillet plutôt qu'une chemise
% de six cents. Le site s'en sert aussi pour tourner les pages du fac-similé
% quand on fait défiler la transcription.
\newcommand{\leaf}[1]{%
  \marginnote{\footnotesize\color{gray!70}\textsf{#1}}%
  \label{leaf:#1}\ignorespaces}

% Illisible : on ne devine pas. Un mot inventé qui se lit comme les autres est
% le pire résultat possible.
\newcommand{\ill}{\textcolor{red!60!black}{[\,\ldots\,]}}

% Lecture douteuse : le mot est donné, et signalé comme incertain.
\newcommand{\uncertain}[1]{\uline{#1}}

% Ajout de l'éditeur — une lettre manquante, un mot sous-entendu.
\newcommand{\add}[1]{\textcolor{blue!50!black}{[#1]}}

% Biffé par Grothendieck lui-même. Ce qu'il a rayé fait partie du document :
% une rature dit souvent où la pensée a changé de direction.
\newcommand{\struck}[1]{\sout{#1}}

% Note du transcripteur, sur l'état matériel du feuillet.
\newcommand{\note}[1]{\par\smallskip{\small\color{gray!60!black}\textit{[#1]}}\par\smallskip}

% Note marginale de Grothendieck, distincte du corps du texte.
\newcommand{\marginal}[1]{\par\smallskip\noindent\hspace{1em}%
  {\small\color{gray!60!black}#1}\par\smallskip}

% --- Titre ------------------------------------------------------------------

\AtBeginDocument{%
  \begin{center}
    {\large\bfseries Fonds Alexandre Grothendieck --- cote n\textsuperscript{o}~\grothFolder}\\[2pt]
    {\itshape\grothFoldertitle}\\[4pt]
    {\small feuillets \grothFirst--\grothLast\ (lot \grothBatch)}\\[2pt]
    {\footnotesize\color{gray!60!black}Transcription établie d'après le fac-similé numérisé
     par l'Université de Montpellier.\\
     Les lectures incertaines sont soulignées, les illisibles notées [\,\ldots\,],
     les ajouts entre crochets.}
  \end{center}
  \vspace{1em}}

\endinput


\folder{120}
\batch{1}
\pages{1}{20}
\dating{[à partir de 1973-à partir de 1978]}
\watermark{Édition de démonstration}
\foldertitle{Topologie modérée : notes manuscrites (s.d.)}

\begin{document}

\section*{Topologie modérée}

\note{inscrit à l'encre en haut à droite du feuillet de garde (page 1), qui ne
reçoit pas de numéro de page : un premier mot, biffé et illisible
(\uncertain{Structures}), est remplacé en interligne par « Topologie », et la
finale de « modérées » est corrigée en « modérée » ; le tout est souligné.
Au-dessous, au crayon très pâle et peut-être d'une autre main, « 1978 \ill{} »
et « (\uncertain{Autour de l'Éloge}) »}

\page{2}

\note{p. 1 de sa pagination}

\marginal{(I) Modèles $R$-modérés \struck{et \ill{}}}\note{titre écrit en biais dans la marge
supérieure gauche, le chiffre romain cerclé ; « $R$-mo » est ajouté au-dessus de la ligne}

$\forall n \in \mathbb{N}^*$, $\mathcal{M}_n \subset \mathfrak{P}(R^n)$

\begin{itemize}
  \item[(M1)] $\forall\, u : [1,n] \to [1,m]$, d'où $u^* : R^m \to R^n$, on a
  $u^*(\mathcal{M}_m) \subset \mathcal{M}_n$ i.e.\ $\forall A \in
  \mathcal{M}_m$, on a $u^*(A) \in \mathcal{M}_n$
  \item[(M2)] $\forall\, u : [1,n] \to [1,m]$ \emph{surjectif}, d'où $u^* :
  R^m \hookrightarrow R^n$ injectif, on a $u^{*-1}(\mathcal{M}_n) \subset
  \mathcal{M}_m$ i.e.\ $\forall B \in \mathcal{M}_n$, on a $u^{*-1}(B) \in
  \mathcal{M}_m$\note{à la fin de la ligne la lettre est surchargée ($B$ sur
  $A$, ou l'inverse) et l'exposant de $u$ est peu net}
  \item[(M3)] $\forall\, u : [1,n] \to [1,m]$ d'où $u^* : R^m \to R^n$, et
  pour $A \in \mathcal{M}_m$, $B \in \mathcal{M}_n$, on a
  \[ u^{*-1}(B) \cap A \in \mathcal{M}_m \]
  \item[(M4)] $\forall\, A \in \mathcal{M}_m$, $B \in \mathcal{M}_n$, on a
  (moyennant $R^m \times R^n \simeq R^{m+n}$)
  \[ A \times B \in \mathcal{M}_{m+n} \]
\end{itemize}
\note{les quatre étiquettes (M1)--(M4) sont cerclées ; les chiffres de (M1),
(M2), (M3) sont repassés}
\marginal{\ill{} M1, \ill{} $A \in \mathcal{M}_m \Rightarrow u_*(A) \in
\mathcal{M}_n$}\note{en biais, dans la marge gauche, en face de (M2)--(M3) ;
le $u_*$ est une lecture douteuse}

\emph{Commentaires} (1) M1) permet de définir, pour \struck{tout} $I$ ensemble
fini non vide, $\mathcal{M}_I \subset \mathfrak{P}(R^I)$, et alors on a des
propriétés $\mathcal{M}'1$ à $\mathcal{M}'4$ pour les $\mathcal{M}_I$ et les
$u : I \to J$.

(2) Pour avoir M1, il suffit de le supposer pour les \struck{\ill{}} cas
\begin{itemize}
  \item[a)] $m = n$, $u \in \mathfrak{S}_n$
  \item[b)] $m = n+1$, $u$ injectif (« op.\ face ») donc $u^*$ surjectif
  \item[c)] $n = m+1$, $u$ surjectif \struck{« op.\ dégénérescence »} (« op.\
  dégénérescence ») donc $u^*$ injectif
\end{itemize}
\struck{et non seulement \ill{} les \uncertain{opérations} \ill{}}
\marginal{on \uncertain{peut se borner} seulement \ill{} $R^n \to R^{n+1}$,
$(x_1, \dots, x_n) \mapsto (x_1, \dots, x_n, x_n)$ et $R^{n+1} \to R^n$,
$(x_1, \dots, x_{n+1}) \mapsto (x_1, \dots, x_n)$}\note{en biais dans la marge
gauche, reliée par une accolade aux cas b) et c)}

Pour avoir M2, il suffit de le supposer pour les op.\ de dégénérescence, et
\uncertain{même seulement} pour l'\ill{} \struck{\ill{}} \ill{},
\marginal{$R^n \hookrightarrow R^{n+1}$, $(x_1, \dots, x_n) \mapsto (x_1,
\dots, x_n, x_n)$}\note{cerclé dans la marge gauche ; une première flèche
biffée}
Notons que M1, M2, M4 impliquent que les $\mathfrak{S}_n$ sont stables par
intersection, car de deux
\[ A, B \in \mathfrak{S}_n \Rightarrow A \cap B \in \mathfrak{S}_n \]
(moyennant M1, M2, M4)\note{le même sigle que le groupe symétrique
$\mathfrak{S}_n$ du cas a) ; on attendrait ici $\mathcal{M}_n$. Transcrit tel
qu'écrit}, d'où résulte que pour avoir M3, il suffit de le prouver pour $u$
surjectif, et même seulement pour $u^*$ l'op.\ standard $R^{n+1} \to R^n$,
$(x_1, \dots, x_{n+1}) \mapsto (x_1, \dots, x_n)$.

\page{3}
En résumé on voit que M1, M2, M3 peuvent être remplacés par
\begin{itemize}
  \item[a)] \struck{\ill{}} $\forall n \in \mathbb{N}^*$, $\sigma \in
  \mathfrak{S}_n$, $A \in \mathcal{M}_n$, on a $\sigma^*(A) \in \mathcal{M}_n$
  \item[b)] $\forall n \in \mathbb{N}^*$, d'où $\alpha_n : [1, \dots, n+1] \to
  [1, n]$ (identité sur $[1,n]$, $n+1 \mapsto n$), d'où \struck{$R$}
  $\alpha_n^* : R^n \hookrightarrow R^{n+1}$, $(x_1, \dots, x_n) \mapsto (x_1,
  \dots, x_n, x_n)$, on a
  \item[b$_1$)] $\forall A \in \mathcal{M}_n$, $\alpha_n^*(A) \in
    \mathcal{M}_{n+1}$,
    \item[b$_2$)] $\forall B \in \mathcal{M}_{n+1}$, $\alpha_n^{*-1}(B) \in
    \mathcal{M}_n$
  \item[c)] $\forall n \in \mathbb{N}^*$, d'où $\beta_n : [1,n]
  \hookrightarrow [1, n+1]$ inclusion can., d'où $\beta_n^* : R^{n+1} \to
  R^n$, $(x_1, \dots, x_{n+1}) \mapsto (x_1, \dots, x_n)$, on a
  \item[c$_1$)] $\forall A \in \mathcal{M}_{n+1}$, $B \in \mathcal{M}_n$, on
    a $A \cap \beta_n^{*-1}(B) \in \mathcal{M}_{n+1}$,
    \item[c$_2$)] $\forall A \in \mathcal{M}_{n+1}$, $\beta_n^*(A) \in
    \mathcal{M}_n$
\end{itemize}
\note{en a), une première écriture est biffée d'un gribouillis ; la lettre
de $\sigma^*(A)$ est surchargée. En b$_1$), l'exposant de $\alpha_n$ est mal
formé ; on lit une image directe}

(3) M1 implique que si un des $\mathcal{M}_n$ est vide, tous le sont --- et
alors les axiomes M1 à M4 sont satisfaits trivialement. On va exclure ce cas
trivial, i.e.\ supposer
\[ (*) \qquad \mathcal{M}_1 \neq \emptyset . \]
M1 implique que si un des $\mathcal{M}_n$ contient $\emptyset$, tous le
contiennent. Si ce n'est pas le cas, alors chaque $\mathcal{M}_n$ (étant
stable par intersection de deux ens.) est une base de filtre sur $R^n$. On
exclut ce cas, i.e.\ on suppose la relation suivante, plus forte que $(*)$
\[ (**) \qquad \emptyset \in \mathcal{M}_1 \quad (\text{d'où } \emptyset \in
\mathcal{M}_n \ \forall n \in \mathbb{N}) \]

\emph{NB} Si cette propriété n'était \uncertain{vérifiée}, alors posant

\page{4}
\note{p. 2 de sa pagination}
\[ \mathcal{M}'_n = \mathcal{M}_n \cup \{\emptyset\} \]
les $\mathcal{M}'_n$ satisfont encore les axiomes \ldots

\struck{(\ill{}) \ill{} \ill{}}\note{un départ biffé, sous un numéro
cerclé lui-même surchargé}
De même, M1 implique que si on a, pour un $n$, \struck{$\mathcal{M}$}
$\mathcal{M}_n = \{\emptyset\}$, alors c'est vrai aussi pour tout $n$ ---
cas trivial (qui satisfait nos axiomes) que je vais exclure. Donc je suppose
aussi
\[ (***) \qquad \mathcal{M}_1 \neq \{\emptyset\} \quad \text{i.e.} \quad
\exists A \in \mathcal{M}_1,\ A \neq \emptyset \]
Ceci étant, si on pose pour \uncertain{$I \neq \emptyset$}
\[ \mathcal{M}_\emptyset = \mathfrak{P}(R^\emptyset) = \{\emptyset,
\{e\}\} \]
($R^\emptyset$ ensemble ponctuel $\{e\}$), les propriétés $\mathcal{M}'1$ à
$\mathcal{M}'4$ (cf.\ (1)) restent satisfaites \ill{} \struck{\ill{}}
$I, J \neq \emptyset$ \ldots\note{ajout interlinéaire entouré d'un trait qui
le rattache à « Ceci étant » ; « $I \neq \emptyset$ » est écrit au-dessus de
« pour », son symbole peut-être biffé}

(4) Soit \struck{$x \in R$},
\[ x = (x_1, \dots, x_n) \in R^n \]
alors il résulte de M1 et \uncertain{M4} que
\[ \{x\} \in R^n \iff \forall i \in [1,n], \text{ on a } \{x_i\} \in
\mathcal{M}_1 . \]
\note{à gauche, « $\{x\} \in R^n$ » : lire $\{x\} \in \mathcal{M}_n$}
Soit
\[ R_0 = \{x \in R \mid \{x\} \in \mathcal{M}_1\} \qquad \text{donc } R_0
\subset R , \]
alors pour tout ens.\ fini $I$, on a
\[ R_0^I = \{x \in R^I \mid \{x\} \in \mathcal{M}_I\} \qquad (R_0^I \subset
R^I) \]
Notons que si \struck{$R_0 \neq \emptyset$} $\operatorname{card} R_0
\geqslant 2$, alors on conclut de la stabilité de $\mathcal{M}_1$ par
intersection que $\emptyset \in \mathcal{M}_1$, donc on satisfait à (3).

(5) Exemples (à part $\mathcal{M}_n = \emptyset$ $\forall n$, et
$\mathcal{M}_n = \{\emptyset\}$ $\forall n$)

\page{5}
\begin{itemize}
  \item[a)] $\mathcal{M}_n$ = parties de cardinal $\leqslant 1$ de $R^n$
  \item[b)] $\mathcal{M}_n$ = parties finies de \ill{} $R^n$
  \item[c)] $\mathcal{M}_n$ = toutes les \ill{} parties de $R^n$
\end{itemize}
\marginal{$\to$ \ill{} les axiomes \ill{} b) c) \uncertain{du NB} (2), \ill{}
l'exclusion de (2)}\note{à droite de la liste, relié par une flèche au cas
a) ; lecture très incertaine}

Aucun de ces \ill{} \ill{} pour nous !\note{la ligne est coupée par un long
ajout interlinéaire écrit en biais, en quatre ou cinq lignes serrées, dont on
ne lit que des fragments : « \ill{} \uncertain{exemples} \ill{} des
\ill{} \uncertain{discrets} \ill{} » ; il n'est pas transcrit}

\subsection*{(II) Espaces modérés}
\note{titre écrit en biais dans la marge gauche, le chiffre romain cerclé,
« Espaces modérés » souligné}

\marginal{Dans ce qui suit, au lieu de (M1) \ill{}, il suffit l'axiome
\ill{} \struck{\ill{}} $[1,n]$ \ill{} \uncertain{un seulement} \ill{} ; on
n'exclut \uncertain{pas} $I$ vide}\note{en biais dans la marge gauche, sous
le titre}

\uncertain{Déf} \ill{} : partie modérée de $R^I$, points modérés de $R$, de
$R^I$.

Soit $A \in \mathcal{M}_I$, $B \in \mathcal{M}_J$, $f : A \to B$, d'où
\[ \Gamma_f \subset A \times B \subset R^I \times R^J \simeq R^{I \sqcup J} .
\]
On dit que $f$ est \emph{modérée} si $\Gamma_f \in \mathcal{M}_{I \sqcup J}$.
\note{il écrit l'indice $I \cup J$ ; c'est la réunion disjointe qui est
entendue}

\marginal{via M2}\emph{Proposition}
\begin{itemize}
  \item[a)] $\mathrm{id}_A$ (si $A \in \mathcal{M}_I$) est modérée
  \item[b)] $f : A \to B$ et $g : B \to C$ modérées ($A \in \mathcal{M}_I$, $B
  \in \mathcal{M}_J$, $C \in \mathcal{M}_K$) $\Rightarrow$ $gf$ modérée
  \item[c)] $f : A \to B$ ($A \in \mathcal{M}_I$, $B \in \mathcal{M}_J$)
  modérée et bijective $\Rightarrow$ $f^{-1}$ modérée
\end{itemize}
(on \uncertain{précise} $(I, A \in \mathcal{M}_I)$)

D'où catégorie des \emph{modèles modérés} ($\mathrm{Mod}_{\mathcal{M}_*}$), et
un foncteur \emph{conservatif fidèle}
\[ \mathrm{Mod}_{\mathcal{M}_*} \longrightarrow (\mathrm{Ens}) \qquad (I, A)
\mapsto A . \]
\marginal{on n'exclut \uncertain{pas} les \uncertain{modèles} \ill{}
\uncertain{dans lesquels} $R_0 = \{e\}$ ! \struck{\ill{}} \ill{}
\uncertain{évident} \ill{} \ldots}\note{en biais dans la marge gauche}
On en déduit la catégorie des ensembles ($\mathcal{M}_*$-)modérés (ou
« espaces »), équivalente à $\mathrm{Mod}_{\mathcal{M}_*}$, mais avec un
foncteur
\[ \mathrm{Esp}_{\mathcal{M}_*} \longrightarrow (\mathrm{Ens}) \qquad X
\mapsto |X| \]
qui est fidèle, conservatif et de plus « transportable », i.e.

\page{6}
\note{sa pagination, en haut à droite, est à demi effacée ; la suite 1, 2, 4, 5 des pages 2, 4, 8, 10 fait attendre 3}
pour tout diagramme
\[ \begin{cases} X \\ |X| \xrightarrow[\sim]{\ \alpha\ } E \end{cases} \]
avec $X \in \operatorname{Ob} \mathrm{Esp}_{\mathcal{M}_*}$, $\exists$ un unique
\struck{$Y$ et} \uncertain{flèche} $f : X \to Y$ de
$\mathrm{Esp}_{\mathcal{M}_*}$, qui relève $\alpha$ (« transport de
structure »).\note{le $X$ du diagramme est écrit sur un $A$}

\emph{Proposition} Dans la catégorie $\mathrm{Esp}_{\mathcal{M}_*}$ des
espaces $\mathcal{M}_*$-modérés, il y a des produits finis, et des
\struck{fibrés de deux objets sur un troisième)} produits (de deux objets), et
le foncteur $X \mapsto |X|$ y commute. Si on suppose
\[ (\mathrm{M}5) \qquad R_0 \neq \emptyset \]
alors il y a \struck{foncteur} un objet final (qui est un ponctuel) et des
\struck{foncteurs} limites finies, et le foncteur $X \mapsto |X|$ y commute.
\note{le passage qui va de « le foncteur $X \mapsto |X|$ » à « limites
finies » est encadré et barré de trois longs traits obliques ; en biais, sur
la droite, un ajout que l'on ne lit que par fragments : « \ill{}
\uncertain{existence} d'un objet final \ill{} des produits \ill{} \ill{}
$X \mapsto |X|$ \ill{} commute \ill{} ». La ligne des produits porte
d'abord « produits \uncertain{fibrés} », repris en « produits (de deux
objets) »}

\emph{Proposition} \struck{\ill{}} Si pour un $n \in \mathbb{N}^*$ on a
\struck{\ill{}} $R^n \in \mathcal{M}_n$, \uncertain{il en est ainsi} pour
tout $I$ fini \uncertain{dans} $R^I \in \mathcal{M}_I$. Considérons l'axiome
(\emph{facultatif})
\[ (\mathrm{M}0) \qquad R \in \mathcal{M}_1 \qquad (\Rightarrow \forall I \
(\text{ens.\ fini}),\ R^I \in \mathcal{M}_I) \]
\note{cette fin de page est recouverte, dans la marge gauche et entre les
lignes, par un long ajout écrit en biais, en une dizaine de lignes, où l'on
distingue « Quand on \ill{} (M0) », « (M5) », « des $A \in \mathcal{M}_1$ »,
« $R$ \ill{} $\Sigma$ », « (M11) » et « \ill{} des \ill{} » ; il n'est pas
transcrit. Devant (M0), une première étiquette est biffée}

\page{7}
Alors $R$ est muni d'une structure modérée canonique (car \struck{\ill{}} que
modèle $\subset R^1 = R$), i.e.
\[ R \in \operatorname{Ob} \Sigma \qquad (\Sigma = \mathrm{Esp}_{\mathcal{M}_*})
\]
et tout objet de $\Sigma$ est isomorphe à un sous-objet d'un $R^n$ (produit
dans $\Sigma$). On reconnaît les $\mathcal{M}_n$ comme les ensembles des
sous-objets (dans $\Sigma$) de l'objet $R^n$. On trouve que la donnée d'un
$\mathcal{M}_*$ est équivalente à la donnée de la \uncertain{structure}
\note{ajout interlinéaire : « \ill{}, les $R^n$ correspondants \ill{} »}
: la donnée de la
\[ (\Sigma, R, X \mapsto |X|) \]
(à la \uncertain{notion} \uncertain{d'isom.\ près}) \ldots\ \ill{}
\begin{itemize}
  \item[a)] la catégorie $\Sigma$ admet des $\varprojlim$ finies
  \item[b)] le foncteur $X \mapsto |X|$, $\Sigma \to (\mathrm{Ens})$, est
  \ill{} \ill{} \ill{}
  \item[c)] l'objet $R \in \operatorname{Ob} \Sigma$ est tel que tout objet
  \ill{} de $R$ \ill{} \ill{} un $R^n$.
\end{itemize}
\note{la fin de la page est très chargée : une note en biais dans la marge
gauche, où l'on lit « Intr\uncertain{oduire} les parties modérées
\uncertain{des} espaces modérés, \ill{} des \ill{} $\Sigma$ », « les
parties \ill{} $f(A)$, $f^{-1}(B)$ », « \ill{} produits cartésiens » ; et
une autre, en biais, en bas à droite, où l'on lit « (M1)$_0$ \ill{} \ill{}
\ill{} » ; elles ne sont pas transcrites}

Parlons aussi des \uncertain{opérations} \ill{} \ill{} \emph{modérées}
\uncertain{d'espaces modérés} \ill{} \emph{dans des $R^I$}.

\subsection*{(III) Limites inductives finies d'espaces modérés}
\note{le chiffre romain est cerclé ; titre souligné, « inductives » ajouté
en interligne}
À partir \ill{} \ill{} \ill{} M1 \ill{} \ill{} \ill{} \ill{} \ill{}
\uncertain{intéressant} : le \uncertain{problème} de quotients des espaces
modérés, et des \emph{limites inductives} finies dans $(\mathrm{Esp}_{\mathcal{M}_*})$.
\note{lecture très incertaine de toute cette fin de page}

\page{8}
\note{p. 4 de sa pagination}
\marginal{la donnée \uncertain{modérée} \ill{} précise le \ill{}, qui \ill{}
\uncertain{facilement} \ill{} \ldots}\note{en biais dans la marge supérieure
gauche}
\emph{Proposition} Si \struck{$R_0 \neq \emptyset$} $\operatorname{card} R_0
\geqslant 2$, \struck{alors dans $\Sigma$ (\ill{} $\Sigma$
($\Leftrightarrow R \in \mathcal{M}_1$)) existe un \ill{}
$(\mathrm{Hom}_\Sigma(e, R) \simeq \dots)$ de l'objet \ill{}
$\mathrm{Hom}_\Sigma(R^0, R) \neq \emptyset$ (avec $R^0 = \{e\}$)} alors
dans $\Sigma$ \struck{il y a des \ill{} \ill{}} \uncertain{limites}
\uncertain{projectives} finies (\uncertain{moyennant} (M5)) et $X \mapsto
|X|$ y commute.\note{les quatre premières lignes de l'énoncé sont biffées
et surchargées ; seule la reprise finale est sûre dans son sens}

\emph{Corollaire} Sur tout ens.\ fini il y a une structure modérée et une
seule, pour laquelle les \uncertain{parties} \uncertain{modérées} soient
\ill{} (str.\ mod.\ discrète).\note{la ligne sous « modérée et une seule »
porte une surcharge interlinéaire peu lisible}
Pour \uncertain{ne} regarder \ill{} \ill{} \ill{} \ill{} construire
\[ X_1 \sqcup_Z X_2 , \]
sans supposer $Z$ fini. Notons d'abord ceci :

\emph{Proposition} Soit $I$ un \struck{type de} diagramme fini, et
$(X_i)_{i \in I}$ un diagramme de type $I$ dans $\Sigma =
\mathrm{Esp}_{\mathcal{M}_*}$, i.e.\ d'espaces modérés. Soit $X$ la limite
inductive \emph{dans} $(\mathrm{Ens})$ (quotient de $\coprod X_i$ \uncertain{par} une
certaine \uncertain{relation}). Supposons qu'il existe sur $X$ une structure
modérée, telle que les \uncertain{applications} can.\ $g_i : X_i \to X$
soient modérées. Alors $X$, muni de cette structure, est une $\varinjlim$ des
$X_i$.\note{« $\varinjlim$ » : il écrit « un $\varinjlim_I$ », l'indice
sous la flèche}

\page{9}
Axiome supplémentaire
\[ (\mathrm{M}5) \qquad \forall n,\ \mathcal{M}_n \text{ stable par réunions
finies} \]
($\Leftrightarrow$ $\emptyset \in \mathcal{M}_n$, et $A, B \in \mathcal{M}_n
\Rightarrow A \cup B \in \mathcal{M}_n$)\note{l'étiquette (M5) a déjà servi,
page 6, pour $R_0 \neq \emptyset$ ; elle est reprise ici telle quelle}

\emph{Proposition} Parties modérées d'un esp.\ mod.\ stable par réunions
finies.

\emph{Proposition} Sur $\emptyset$ il y a une structure modérée unique, qui
en fait objet initial de $\mathrm{Esp}_{\mathcal{M}_*}$.

\emph{Proposition} Soient $X$ un espace modéré, $(X_i)_{i \in I}$ une famille
finie de parties modérées de $X$. Alors $\bigcup X_i$ est une partie modérée
de $X$. Si $X = \bigcup_i X_i$, alors $X$ est somme amalgamée des $X_i$ suivant
les $X_i \cap X_j$ --- une application $f : X \to Y$ entre espaces modérés
est modérée ssi ses restrictions aux $X_i\,$ le sont.

\emph{Question} Donnons-nous

\begin{tikzcd}[column sep=small]
 & Z \arrow[dl, hook'] \arrow[dr, hook] & \\
X_1 & & X_2
\end{tikzcd}

dans $\mathrm{Esp}_{\mathcal{M}_*}$, existe-t-il une somme amalgamée ? Si $Y =
\emptyset$, c'est la question de l'existence des sommes finies dans
$(\mathrm{Esp}_{\mathcal{M}_*})$.\note{le $Z$ du diagramme est écrit sur une
autre lettre, sans doute $Y$, qui reste dans « Si $Y = \emptyset$ »}

\page{10}
\note{p. 5 de sa pagination}
Pour que la somme amalgamée existe dans \struck{Esp} $\Sigma$ et que $X
\mapsto |X|$ y commute, il faut et suffit qu'il existe une famille finie de
\uncertain{paires de} fonctions
\[ \begin{cases} f_\alpha : X_1 \to R \\ g_\alpha : X_2 \to R \end{cases} \]
telles que l'on ait
\begin{itemize}
  \item[a)] $f_\alpha(i_1(z)) = g_\alpha(i_2(z))$ \quad $\forall \alpha \in A$,
  $z \in Z$
  \item[b)] \struck{Si $x, y$ les $f_\alpha$ \ill{} ($f_\alpha$) :} les
  $f_\alpha$ ($g_\alpha$) séparent les pts de $X_1$ ($X_2$) i.e.
  \[ X_1 \to R^I \qquad (X_2 \to R^A) \]
  sont injectifs
  \item[c)] Si $x_1 \in X_1 \setminus i_1(Z)$, $x_2 \in X_2 \setminus
  i_2(Z)$, $\exists \alpha$ avec $f_\alpha(x_1) \neq g_\alpha(x_2)$
\end{itemize}
\note{l'ensemble d'indices est une lettre $A$ cursive, que l'on transcrit
$A$ ; en b), le premier exposant ressemble à $I$, transcrit tel quel. En a),
le premier indice $i_1$ est peu net. En c), une première lettre après
$x_1$ est biffée}

\emph{NB} On peut généraliser, sous les conditions de la proposition, il
existe la structure modérée \uncertain{quotient} sur $X$, ssi $\exists$ une
famille finie $(F_\alpha)$ \ill{} de systèmes \uncertain{compatibles}
\[ (f_{\alpha i})_{i \in I} \qquad f_{\alpha i} : X_i \to R \text{
modérée} , \]
\uncertain{définissant} des
\[ f_\alpha : X \to R \]
telles que les $f_\alpha$ séparent les pts de $X$, i.e.\ $f = (f_\alpha) : X
\hookrightarrow R^A$

\page{11}
En effet, si $Y$ est un espace modéré, une \struck{\uncertain{bijection}}
\ill{} des \uncertain{applications} \uncertain{inductif} $X_i$ (dans $\Sigma$)
dans $Y$, $(f_i : X_i \to Y)_{i \in I}$, c'est-à-dire la donnée d'une
application $f : X \to Y$, telle que les $f \circ g_i$ soient modérées. Il
est évident : prouvons que cela implique que $f$ est modérée, i.e.\ que son
graphe $\Gamma_f \subset X \times Y$ est une partie modérée. Or $X$ étant la
réunion des images des $X_i$, il en est que $\Gamma_f$ est la réunion des
images des $\Gamma_{f_i} \subset X_i \times Y$ par
\[ X_i \times Y \xrightarrow{\ g_i \times \mathrm{id}_Y\ } X \times Y . \]
Mais ces images sont modérées, et donc leur réunion aussi.
\note{le début de la page est d'une lecture incertaine ; la famille
$(f_i : X_i \to Y)_{i \in I}$ est écrite au-dessus de la ligne. Dans
$X_i \times Y \to X \times Y$, un premier facteur est biffé}

\emph{Cor} Soit $f : X \to Y$ application modérée, $\mathcal{R}_f$ la relation
d'équivalence définie par $f$ (qui est donc une partie modérée de $X \times
X$, et est une $\Sigma$-relation d'équivalence). Alors $X / \mathcal{R}_f$
existe dans $\Sigma = \mathrm{Esp}_{\mathcal{M}_*}$, et est can.\ isomorphe
à $f(X)$ (avec structure induite).\note{« can. » est ajouté au-dessus de la
ligne}

\emph{Cor.\ 2} Considérons

\begin{tikzcd}[column sep=small]
 & Z \arrow[dl, hook'] \arrow[dr, hook] & \\
X_1 & & X_2
\end{tikzcd}

dans $\Sigma$.\note{le diagramme, écrit en ligne, est traversé de quatre
traits obliques, qui semblent annuler la ligne ; la suite est page 12}

\page{12}
\note{p. 6 de sa pagination}
On va introduire
\[ (\mathrm{M}6) \qquad \begin{cases}
\forall\, Z, X \in \mathcal{M}_n,\ Z \subset X,\ \exists\,
(f_\alpha)_{\alpha \in A} \text{ famille finie de fonctions modérées } X \to
R \\
\text{et des } (e_\alpha)_{\alpha \in A},\ e_\alpha \in R_0, \text{ telles
que } Z = \bigcap f_\alpha^{-1}(e_\alpha) . \\
\text{De plus OPS que les } f_\alpha \text{ séparent les points de } X
\smallsetminus Z
\end{cases} \]
\note{devant l'accolade, un premier $\exists$ est biffé ; « $Z, X$ » est
écrit sur une première lettre}
\struck{Il en résulte que les lim.\ finies des diagrammes (cont.)}

\struck{(\ill{})}

\begin{tikzcd}[column sep=small, row sep=small]
X \arrow[r, hook] & R^n \times R^A \\
Z \arrow[u, hook] \arrow[r, hook] & R^n \times \lbrace e \rbrace
\arrow[u, hook]
\end{tikzcd}

$e = (e_\alpha)_{\alpha \in A}$

\struck{Il s'ensuit ceci : $Z_1 \subset X_1$ et $Z_2 \subset X_2$}
\note{tout le bloc entre la ligne « Il en résulte \ldots » et « $Z_2 \subset
X_2$ » est encadré de deux traits horizontaux et barré de quatre traits
obliques ; le diagramme est pris dans la biffure, on le donne néanmoins}

\uncertain{Notons} que \emph{signifiant} en fait que le quotient $X/Z$
(contraction de $Z$ en un point) existe dans $\Sigma$ en un sens strict
(i.e.\ $X \mapsto |X|$ y commute). De plus, on supposera
\[ (\mathrm{M}7) \qquad \begin{cases} \forall\, Z, X \in \mathcal{M}_n,\ Z \subset X, \text{
et toute } f : Z \to R \text{ modérée,} \\ \exists\, g : X \to R \text{
modérée qui prolonge } f \end{cases} \]

\emph{Proposition} Moyennant M5, M6, M7, les sommes amalgamées
\[ X_1 \sqcup_Z X_2 \qquad (Z \hookrightarrow X_1,\ Z \hookrightarrow X_2) \]
existent dans $\mathrm{Esp}_{\mathcal{M}_*}$, et commutent au foncteur
oubli\ldots
\marginal{(M6) (M7) implique l'existence \ill{} \ill{} \ill{} \ill{}
\uncertain{sommes} \uncertain{amalgamées}}\note{en biais dans la marge
gauche, en plusieurs lignes, en face du diagramme et de la Proposition ;
lecture très incertaine}

\page{13}
Soit $Z \hookrightarrow R^I$ un plongement modéré. On peut (par M7) le
prolonger en
\[ X_i \longrightarrow R^I \qquad (i \in \lbrace 1, 2 \rbrace) \]
Soit
\[ X_i \longrightarrow R^{A_i} \qquad i \in \lbrace 1, 2 \rbrace \]
comme dans M6 (qui \uncertain{définit} un \struck{\ill{}} plongement modéré
$X_i / Z_i \hookrightarrow R^{A_i}$, $Z_i = u_i(Z)$), \struck{on} on a donc
un
\[ X_1 \to R^{I \sqcup A_1} \]
qui est un plongement et s'insère dans un diagramme cartésien

\begin{tikzcd}[column sep=small, row sep=small]
X_1 \arrow[r, hook] & R^{I \sqcup A_1} = R^I \times R^{A_1} \\
Z \arrow[u, hook] \arrow[r, hook] & R^I \times \lbrace e_1 \rbrace
\arrow[u, hook]
\end{tikzcd}

$e_1 \in R_0^{A_1}$\note{dans ce diagramme, la flèche verticale de gauche
est tracée vers le bas, de $X_1$ vers $Z$ ; on la donne dans le sens de
l'inclusion, qui est celui de la flèche de droite}

et de même pour $X_2$, \struck{\ill{}} ce qui permet de considérer comme
\struck{dans} s'insérant dans
\note{suit un premier diagramme, sur deux lignes, barré de six traits
obliques : $X_1 \hookrightarrow R^I \times R^{A_1} \times \lbrace e_2
\rbrace \to R^I \times R^{A_1} \times R^{A_2}$ au-dessus de $Z \to R^I
\times \lbrace e_1 \rbrace \times \lbrace e_2 \rbrace \hookrightarrow R^I
\times R^{A_1} \times R^{A_2}$ ; il est repris, corrigé, dans le
diagramme (*) qui suit}

\begin{tikzcd}[column sep=small, row sep=small, nodes={font=\scriptsize}]
X_1 \arrow[r, hook] & R^I \times R^{A_1} \times \lbrace e_2 \rbrace &
R^I \times R^{A_1} \times R^{A_2} \arrow[l] \\
 & R^I \times \lbrace e_1 \rbrace \times \lbrace e_2 \rbrace \arrow[u, hook]
\arrow[r, hook] & R^I \times \lbrace e_1 \rbrace \times R^{A_2} \arrow[u,
hook] \\
Z \arrow[uu, hook] \arrow[rr, hook] & & X_2 \arrow[u, hook]
\end{tikzcd}

\note{diagramme marqué (*) dans la marge gauche. La flèche de la première
ligne, de $R^I \times R^{A_1} \times R^{A_2}$ vers $R^I \times R^{A_1}
\times \lbrace e_2 \rbrace$, est tracée de droite à gauche, sans crochet ;
les deux crochets de la colonne de droite sont notés par de petites
marques et leur sens est incertain ; le $Z$ du coin inférieur est écrit sur
une autre lettre}

\page{14}
\note{p. 7 de sa pagination}
\marginal{M5 \ill{} \ill{} mais \ill{} M6}\note{dans la marge supérieure
gauche, en biais}

\subsection*{(IV) Ind-structures modérées strictes}
\note{le chiffre romain est cerclé ; « espaces modérés » est écrit au-dessus
de « structures modérées », et en haut à droite, entre parenthèses, « (ou
Ind-espaces modérés) ». Le préfixe s'écrit ici, trois fois, avec une
majuscule nette : « Ind- »}

\ill{} : c'est : la donnée d'un ens.\ $E$, et un ens.\ de parties
$\mathcal{M}_E$ sur $E$ (les parties modérées), stable par réunions finies,
(recouvrant $E$,) avec une structure modérée sur chacune, ces structures
s'induisant mutuellement, et
\[ A \in \mathcal{M}, \ B \in \mathcal{M}_A \Rightarrow B \in \mathcal{M} .
\]
\note{« recouvrant $E$ » est ajouté au-dessus de la ligne ; devant
« $\mathcal{M}_E$ », une première lettre est surchargée, de même devant
« $A \in \mathcal{M}$ »}
Il revient au même de se donner la notion d'application modérée d'un $X$
modéré dans $E$, avec \struck{\ill{}} quelques conditions
\uncertain{simples}\ldots\ Si on suppose
\[ (\mathrm{M}8) \qquad R \text{ est réunion de ses parties modérées} \]
(\uncertain{impliqué} par M5)\note{la parenthèse est écrite au-dessus de
la ligne}, on trouve une structure ind-modérée canonique sur $R$, et de
même sur les $R^I$.\note{ici « ind- » en minuscule, lu comme le « Ind- » du
titre}

\marginal{M1 : M5}\emph{Proposition} La catégorie $\mathrm{Esp}_{\mathcal{M}}$
(des Ind-espaces \uncertain{modérés} stricts) admet des $\varprojlim$
finies, et le foncteur oubli y commute. Si $R_0 \neq \emptyset$ (idem dans
$\mathrm{Esp}_{\mathcal{M}}$ il y a sommes \ill{}) il y a des sommes directes
finies dans $\mathrm{Esp}_{\mathcal{M}}$, et $X \mapsto |X|$ y commute.
\note{« M1 : M5 » en marge, en face de la Proposition, comme « M1 : M7 »
page 23. Le sigle de la catégorie porte un indice surchargé,
$\mathcal{M}_*$ ou $\mathcal{M}$. Dans la marge gauche, en biais, une
remarque dont on lit « $\mathrm{Esp}_{\mathcal{M}}$ \ill{} \ill{} \ill{}
stricts \ill{} \ill{} (en principe !) » ; un premier mot devant
« catégorie » et un devant « dans » sont biffés}

\page{15}
Soit $\mathfrak{X}$ un espace ind-modéré. Une partie de $\mathfrak{X}$ est
dite loc.\ modérée si
\[ \forall A \in \mathcal{M}_{\mathfrak{X}}, \text{ on a } A \cap \mathfrak{Y}
\in \mathcal{M}_{\mathfrak{X}} . \]
C'est stable par intersections finies (y compris pour famille vide, i.e.\
$\mathfrak{X} \in \vec{\mathcal{M}}_{\mathfrak{X}}$), par image inverse par
morphismes \struck{loc.} d'espaces ind-modérés.\note{ses lettres gothiques
$\mathfrak{X}$, $\mathfrak{Y}$, $\mathfrak{Z}$ pour les espaces
ind-modérés ; la première est écrite sur une autre lettre. Il note
$\vec{\mathcal{M}}_{\mathfrak{X}}$, une flèche au-dessus et une au-dessous,
l'ensemble des parties loc.\ modérées}

La notion d'application modérée d'un espace modéré $X$ dans un $R^I$, est un
cas particulier de la précédente, quand $\mathcal{M}_{\mathfrak{X}}$
\uncertain{est} satisfait dans $X$ et $R^I$ sont \uncertain{deux} des
espaces ind-modérés.

Moyennant M6, M7, il existe \struck{pas} dans $\mathrm{Esp}_{\mathcal{M}_*}$
\uncertain{les quotients} $\mathfrak{X}/\mathfrak{Y}$ au \uncertain{sens}
strict (i.e.\ \uncertain{qui} commutent à $X \mapsto |X|$) si $\mathfrak{Y}$
est partie loc.\ modérée, et de même
\[ \mathfrak{X}_1 \sqcup_{\mathfrak{Z}} \mathfrak{X}_2 \]
pour $\mathfrak{Z} \hookrightarrow \mathfrak{X}_1$ et $\mathfrak{Z}
\hookrightarrow \mathfrak{X}_2$ des \uncertain{immersions}
\uncertain{modérées} i.e.\ correspondant à des parties \uncertain{loc.}\
modérées

\page{16}
\note{p. 8 de sa pagination}
$\mathfrak{Z}_i \subset \mathfrak{X}_i$ et $\mathfrak{Z} \simeq
\mathfrak{Z}_1$, $\mathfrak{Z} \simeq \mathfrak{Z}_2$).

\emph{NB} On fera attention que \uncertain{contrairement} à ce qui se passe
dans $\mathrm{Esp}_{\mathcal{M}_*}$) un sous-objet d'un objet $\mathfrak{X}$
de $\mathrm{Esp}_{\mathcal{M}_*}$ i.e.\ d'un esp.\ ind-modéré, ne correspond
pas néc.\ à une partie loc.\ modérée\ldots\note{le premier sigle
$\mathrm{Esp}$ est écrit sur un autre mot ; les deux portent peut-être une
autre marque, qui distinguerait la catégorie des espaces ind-modérés}

\subsection*{(V) Relations avec topologies}
\note{le chiffre romain est cerclé, le titre souligné}
Supposons $R$ muni d'une topologie, donc les $R^I$ et les « modèles »
modérés \uncertain{aussi}. On voudrait que les applications modérées soient
continues --- c'est OK si les parties modérées sont compactes (donc fermées,
si $R$ séparé), i.e.\ si
\struck{Cas où les fonctions $X \mapsto$}
\[ (\mathrm{M}9) \qquad \begin{cases} R \text{ séparé} \\ \text{les parties
modérées de } R^n \text{ sont compactes} \end{cases} \]
\note{« Supposons donc » est écrit au-dessus de la ligne biffée ;
l'étiquette (M9) est cerclée}
Alors le foncteur $X \mapsto |X|$ se factorise par la cat.\ des espaces
topologiques
\[ \mathrm{Esp}_{\mathcal{M}_*} \longrightarrow (\text{Esp.\ top.\ cpts}) \]
\note{dans la parenthèse, un premier « Esp » est surchargé}

\page{17}
et ce foncteur est encore fidèle conservatif transportable (\ill{}
\uncertain{égalité}). \struck{\ill{}} On a une réciproque quand $R \in
\mathcal{M}_1$ (donc $R$ cpct) en termes d'un tel foncteur et d'un objet $R
\in \operatorname{Ob} \Sigma$, tel que tout objet de $\Sigma$ se plonge dans
un $R^n$.\note{la parenthèse après « transportable » est ajoutée au-dessus
de la ligne ; devant $\mathcal{M}_1$, une première lettre est surchargée}

Tout espace ind-modéré sera muni de la top.\ limite inductive. Elle induit
sur chaque partie modérée sa topologie propre, quand on suppose que l'espace
est ind-\uncertain{dénombrable}, i.e.\ $\varinjlim$ de \emph{suites}
croissantes d'espaces modérés.

On voudrait que dans un espace modéré, tout pt ait un syst.\ fond.\ de
voisinages \struck{\ill{}} modérés. Il suffit que \struck{\ill{}} ce soit
vrai dans $R$, donc dans les $R^n$, d'où l'axiome
\[ (\mathrm{M}10) \qquad \text{Dans } R, \text{ tout pt a un syst.\ fond.\
de voisinages} \in \mathcal{M}_1 \]
\note{« l'analogue » est écrit au-dessus de « ce soit vrai »}

\page{18}
\note{p. 9 de sa pagination}
On sait alors que c'est une la notion d'application $X \to Y$ entre espaces
modérés, qui est \emph{modérée au voisinage d'un pt}, et on trouve\ldots

\emph{Prop} Pour qu'une application d'un espace modéré $X$ dans un espace
ind-modéré soit modérée, il faut et s.\ qu'elle le soit au voisinage de
chaque pt.

Donc les applications modérées de $X$ dans $Y$ forment les sections d'un
faisceau --- le faisceau des applications modérées de $X$ dans $Y$. La
structure modérée de $X$ est connue quand on connaît l'ens.\ des
applications modérées de $X$ dans $R$ (cela résulte dans le contexte
purement ensembliste, avec (M1)$_0$ (M2) (M3) (M4) seulement), donc aussi
quand on connaît le faisceau des germes de fonctions modérées.
\note{« modérée » est ajouté au-dessus de « structure » ; « (M1)$_0$ » :
l'indice $0$ est net, cf.\ la note en biais de la page 7}

\page{19}
\emph{Espace loc.\ modéré} Esp.\ \struck{loc.\ cpct} ind-modéré dont la
top.\ est \struck{loc.\ cpcte} séparée, et tout point admet un
\emph{voisinage} qui est \struck{\ill{}} modéré (donc l'espace est loc.\
cpct), et tout pt admet un syst.\ fond.\ de voisinages modérés, et les
parties compactes sont les parties contenues dans une partie modérée --- qui
forment donc un ens.\ filtrant cofinal : celui des parties compactes.
\note{« ind-modéré » est écrit au-dessus de « loc.\ cpct » biffé}

Par exemple, par (M10) $R$ est un espace loc.\ modéré, donc les $R^I$ aussi.

Les espaces loc.\ modérés \struck{$\mathrm{Esp}_{\mathcal{M}_*}$} forment une
sous-catégorie pleine $\mathrm{Esplo}_{\mathcal{M}_*}$ de
$\mathrm{Esp}_{\mathcal{M}_*}$, stable par $\varprojlim$ finies.\note{il
écrit « Esploc », « loc » en indice partiel ; le second
$\mathrm{Esp}_{\mathcal{M}_*}$ est souligné, sans doute pour la catégorie
des espaces ind-modérés}

\marginal{Dire que toute \ill{} \uncertain{est contenue} \ill{} d'un \uncertain{Esp} \struck{\ill{}} \ill{} modérée}\note{en biais,
dans le coin inférieur gauche ; au-dessous, un mot isolé biffé}

\page{20}
\note{p. 10 de sa pagination}

\subsection*{(VI) Relation avec structure de corps}
\note{le chiffre romain est cerclé ; le titre est écrit en biais dans la
marge gauche, souligné, « de corps » au-dessus de la ligne}

Oublions les structures topologiques, et supposons par contre que $R$ est
une \emph{structure de corps}. On supposera alors, \struck{\ill{}} en plus de
(M1) et (M4) (minimum pour construire $\mathrm{Esp}_{\mathcal{M}_*}$ avec ses
$\varprojlim$ finies) l'axiome
\struck{(M11) Pour toute partie modérée $A$ de $R^1$, l'application $a
\mapsto -a$ de $A$ dans $R$ est modérée (i.e.\ l'ens.\ $\lbrace (a, -a) \in
R^2 \mid a \in A \rbrace$ est modérée}
\note{ce premier énoncé de (M11) est barré de sept traits obliques ; un
trait le relie à la reprise qui suit. Après « (M1) », un indice est
surchargé}
\[ (\mathrm{M}11) \qquad \begin{cases}
\text{a) } 0 \neq 1 \text{ sont } \in R_0 \text{ i.e.\ } \lbrace 0 \rbrace,
\lbrace 1 \rbrace \in \mathcal{M}_1 \\
\text{b) } \forall A \in \mathcal{M}_2, \text{ les fonctions } (x,y) \mapsto
x - y, \ (x,y) \mapsto xy \\ \qquad \text{sur } A \text{ sont modérées, i.e.\ les} \\
\qquad \lbrace (x, y, x-y) \mid (x,y) \in A \rbrace \text{ et } \lbrace (x,
y, xy) \mid (x,y) \in A \rbrace \\ \qquad \text{sont modérées dans } R^3 \\
\text{c) } \forall A \subset R^*,\ A \in \mathcal{M}_1,\ \exists B \in
\mathcal{M}_1 \text{ avec } A^{-1} \subset B
\end{cases} \]
\note{l'étiquette (M11) est encadrée. Entre b) et c), une première version
de c), « $\forall A \subset R^*$, $A \in \mathcal{M}_1$, $\exists B \in
\mathcal{M}_1$ tel que $A^{-1} \subset B$ », est biffée ; en c), « $A^{-1}$
modérée » est écrit puis repris, et en marge droite, deux fois, « (OPS $B =
A^{-1}$ !) », la première fois biffée}
\struck{Cela ex\ill{}} (les parties a) b)) équivaut à la validité de la
\ldots
\marginal{NB si \ill{} $A \in \mathcal{M}_2$ \ill{} (M11) pour $A = R^2$}
\note{en biais dans la marge gauche, en face de (M11)}

\emph{Proposition} Pour un espace modéré $X$, les fonctions modérées $X \to
R$ forment un \struck{et $=$ une sous-$R_0$-} \emph{sous-anneau} de l'anneau
des fonctions \ill{} sur $X$. En particulier $R_0$ est un
sous-anneau de $R$, et les \ldots
\marginal{\ill{} \uncertain{particulier} \ill{} \ill{} $R_0$
\uncertain{intègre}}\note{verticalement dans la marge gauche, en bas. La
phrase « En particulier \ldots » s'arrête au bas du feuillet ; la page 21
(lot 2) s'ouvre sur un signe isolé qui n'en est pas visiblement la suite}

\end{document}
